\section{Finite-volume recording and uniform drift}
\label{sec:ex-annulus-recording}

The continuum pole calculation controls a coefficient of the analytically
assembled response. The finite-volume argument needs two additional steps:
an exact determinant identity at each torus size, and a quantitative bound
on the difference between a torus mean and its zero Fourier mode. This
section proves that bound, explains the determinant calculation and
positive-matrix domain, and gives the exact scale order in the final
annulus output.
The local stationary estimates used to feed this construction remain those
listed in the preceding section and its evidence record.

Write $\mathcal H_\sigma(A)$ for the class of periodic functions that are
holomorphic when each imaginary momentum coordinate has modulus less than
$\sigma$ and have supremum norm at most $A$ there. The coefficient of a
translation-invariant two-leg picture is its mixed second derivative; its
finite-volume realization at a conserving pair is $\ell^4$ times that
coefficient. Write $\operatorname{Err}_\ell(\mathcal Y,\mathcal P;\delta)$
when every mixed two-leg coefficient of the finite test $\mathcal P$
differs from this finite-volume realization for $\mathcal Y$ by at most
$\ell^4\delta$.

\subsection{A torus mean differs from the zero mode by an exponential tail}

Let $\ell\geq1$ be an integer and let $f$ be a periodic function of four
real momenta, holomorphic and bounded by $A$ on the complex strip of width
$\sigma>0$. Write its absolutely convergent Fourier series as
\[
 f(\theta)=\sum_{n\in\mathbb Z^4}c_n e^{i n\cdot\theta},
 \qquad |n|_1=\sum_{j=0}^3|n_j|.
\]
The normalized mean over the side-$\ell$ grid is
\[
 \langle f\rangle_\ell
   =\ell^{-4}\sum_{q\in(\mathbb Z/\ell\mathbb Z)^4}
                       f(2\pi q/\ell).
\]
Character orthogonality in each coordinate says
\[
 \ell^{-4}\sum_q e^{2\pi i n\cdot q/\ell}
       =\mathbf 1_{n\in\ell\mathbb Z^4}.
\]
Consequently $\langle f\rangle_\ell- c_0$ is the sum of precisely the
nonzero Fourier coefficients whose indices lie in $\ell\mathbb Z^4$.
Every such index satisfies $|n|_1\geq\ell$. Cauchy's integral formula,
applied after shifting each coordinate contour by $\pm i(\sigma/2)$,
and the geometric-series bound in four coordinates give
\[
 \sum_{n\in\mathbb Z^4}e^{(\sigma/4)|n|_1}|c_n|
 \leq A\left(2\frac{1+\sigma/4}{\sigma/4}\right)^4.
 \tag{RY1}
\]
For completeness, the one-coordinate sum used here is bounded by
\[
 \sum_{k\in\mathbb Z}e^{-h|k|}
 =1+2\sum_{k\geq1}e^{-hk}
 \leq 2\frac{1+h}{h},\qquad h=\sigma/4;
\]
the last inequality follows from $e^h\geq1+h$. The contour shift gives
the factor $e^{-(\sigma/2)|n|_1}$ for an individual coefficient, leaving
the weight $e^{(\sigma/4)|n|_1}$ in (RY1). Therefore
\[
 |\langle f\rangle_\ell-c_0|
 \leq \mathfrak r(\sigma,A)e^{-\sigma\ell/4},\qquad
 \mathfrak r(\sigma,A)
   =A\left(2\frac{1+\sigma/4}{\sigma/4}\right)^4.
 \tag{RY2}
\]
This estimate is for the actual full grid at each $\ell$; no subsequence
selection is involved.

\begin{lemma}[Recorded second-order coefficients]
\label{lem:ex-pole-recording}
Let $F_{ij}(p,q)$ be two-momentum response entries in
$\mathcal H_\sigma(A)$ jointly in $p,q$. Let $\mathcal P_\ell$
be a real finite-volume test and
assume its exact local coefficient identity: for every pair of colored
sites $q,q'$, the mixed second coefficient of $\mathcal P_\ell$ is zero
when $q+q'\ne0$, and otherwise is
\[
 \ell^4\,\langle F_{ij}(\operatorname{mon}(q'),\,\cdot\,)\rangle_\ell.
 \tag{RY3}
\]
Define the internal zero mode and limiting picture by
\[
 \widehat F_{ij,0}(p)=\int_{\mathbb T^4}F_{ij}(p,q)\,dq,
 \qquad
 \mathcal Y_{ij}(p)=\tfrac12\,\widehat F_{ij,0}(p),
\]
where Haar measure on $\mathbb T^4$ is normalized to one. We use the
translation-invariant two-leg convention in which the mixed second
derivative is the polarized coefficient (there is no factorial in the
definition of that coefficient). Thus the finite-volume realization at a
conserving pair is $\ell^4$ times this coefficient, namely
$\ell^4\,2\mathcal Y_{ij}$. Then
$\mathcal Y$ is interpreted as the translation-invariant two-leg picture:
its mixed coefficient at external momenta $p,p'$ is zero unless
$p+p'=0$, and in the conserving case is
$2\mathcal Y_{ij}(p')=\widehat F_{ij,0}(p')$.
$\mathcal Y\in\mathcal H_{\sigma/2}(A)$, and for every side $\ell$,
\[
 \operatorname{Err}_\ell(\mathcal Y,\mathcal P_\ell;
     \mathfrak r(\sigma,A)e^{-\sigma\ell/4}).
 \tag{RY4}
\]
\end{lemma}
\begin{proof}
The normalized Haar integral defining $\widehat F_{ij,0}$ is holomorphic
in $p$ on the width-$\sigma$ strip: on each compact substrip its integrand
and derivatives are uniformly bounded, so differentiation under the
integral is valid. Its norm is at most $A$, hence the half-scaled picture
is bounded by $A$ on the width-$\sigma/2$ strip. For each fixed real $p$,
apply (RY2) to the one-variable function $q\mapsto F_{ij}(p,q)$; its
Fourier zero mode is exactly $\widehat F_{ij,0}(p)$. This gives
\[
 |\langle F_{ij}(p,\cdot)\rangle_\ell
              -\widehat F_{ij,0}(p)|
       \leq \mathfrak r(\sigma,A)e^{-\sigma\ell/4}
\]
on the real external momentum torus. If $q+q'\ne0$, both mixed
coefficients in the definition of $\operatorname{Err}_\ell$ vanish by
momentum conservation. If $q+q'=0$, the second Taylor coefficient of
$\mathcal Y=\frac12\widehat F_{ij,0}$ has the factor $2!=2$, so it equals
$\widehat F_{ij,0}$; subtracting (RY3) and multiplying the mean estimate
by the number $\ell^4$ of sites gives (RY4).
\end{proof}

For a finite side $\ell$, set
\[
 \mathcal Y_{\ell,ij}(p)=\tfrac12\,\ell^{-4}
        \sum_{q\in(\mathbb Z/\ell\mathbb Z)^4}F_{ij}(p,2\pi q/\ell)
\]
at each external grid momentum. By (RY3), its second
finite-volume coefficient is $2\ell^4\mathcal Y_{\ell,ij}(p)$.
The limiting picture also has the exchange and reality symmetries
\[
 \mathcal Y_{ij}(p)=\mathcal Y_{ji}(-p),\qquad
 \overline{\mathcal Y_{ij}(p)}=\mathcal Y_{ij}(-p).
\tag{RY4a}
\]
At each finite side, the logarithmic determinant is a real scalar
function of real field coordinates. Equality of its mixed derivatives
gives exchange of the two external legs, while reality of its Hessian
gives conjugation under Fourier inversion. Thus the finite-volume
normalized pictures have both relations at every allowed grid momentum.
Estimate (RY2) says $\mathcal Y_{\ell,ij}(p)\to\mathcal Y_{ij}(p)$
uniformly over the external grid points as $\ell\to\infty$. Thus their
difference from $\mathcal Y$ tends uniformly to zero on
these grids as the side tends to infinity. The union of the full side
grids is dense in the momentum torus, and $\mathcal Y$ is continuous by
its analytic strip bound. Passing along cofinal sides and then using
density proves (RY4a) everywhere. This supplies the Hessian symmetry and
reality of the half-scaled zero-mode picture; it does not follow from
factorial normalization alone.

The factor $1/2$ in $\mathcal Y$ is essential: the finite-volume mixed
second derivative of $\frac12\widehat F_{ij,0}$ contributes
$2!\,\frac12\widehat F_{ij,0}=\widehat F_{ij,0}$. This is why (RY3)
compares the normalized grid mean to the continuum coefficient and why the
error is the per-site Fourier error multiplied by the site count exactly
once.

\subsection{Determinant identity and positive-matrix domain}

In the application, the pole test is
\[
 \mathcal P_\ell(x)=\tfrac12\log\det P_\ell(x)
                              -\log\det V_\ell(x),
\]
where $P_\ell$ is the positive rank matrix and $V_\ell$ the positive
reference matrix. The finite ticket identifies the zero-, one-, and
two-leg Taylor layers of $P_\ell$ with those of the interacting response,
and those of $V_\ell$ with the reference response. The floor and
positive-regulator hypotheses make both matrices positive definite at
$x=0$; by continuity their real logarithmic determinants are smooth
nearby. Ordinary matrix differentiation gives
\[
 \begin{aligned}
 D\log\det P[H]&=\operatorname{tr}(P^{-1}DP[H]),\\
 D^2\log\det P[H,K]
 &=\operatorname{tr}\!\left(P^{-1}D^2P[H,K]
       -P^{-1}DP[K]P^{-1}DP[H]\right).
 \end{aligned}
 \tag{RY8}
\]
and the same formulas for $V$. Write $A(r)$ for either matrix family,
$B_i(p,r)$ for its first variation (from background momentum $r$ to
$r+p$), and $C_{ij}(p,q,r)$ for its mixed second variation (from $r$ to
$r+p+q$). The shifted fiber block of the mixed derivative of the
logarithmic determinant at the ordered pair of legs is the trace of
\[
 A(r+p+q)^{-1}C_{ij}(p,q,r)
 -A(r+p+q)^{-1}B_i(p,r+q)A(r+q)^{-1}B_j(q,r).
 \tag{RY8a}
\]
This follows by differentiating
$D\log\det A[H]=\operatorname{tr}(A^{-1}DA[H])$ and using
\[D(A^{-1})[K]=-A^{-1}DA[K]A^{-1}.\] There is one ordered product for
the specified ordered pair; exchanging the legs gives the corresponding
other entry, not an extra copy in (RY8a). Sum this expression over internal
momentum and apply the weights in the finite ticket. Formula (RY8a) is the
shifted fiber block of the Hessian. Translation invariance makes the full
trace entry zero unless the external momenta conserve; in the conserving
case summing the fiber blocks gives the normalized torus mean multiplied by
$\ell^4$. This finite-volume determinant
calculation precedes the Fourier tail limit.

The Fourier estimate is entrywise. Any closed-cone passage for finite
reflection matrices and the smooth-word approximation to point
configurations belong to the separate reflection argument.

\subsection{Order of choices in the all-period output}

Here is the exact downstream implication, with its inputs exposed. Fix
$\sigma,\epsilon,\tau>0$ and $C,Z,A\geq0$. Assume a baseline analytic
family $H$ has strip fee $(k,E)$ with $k>0$, $E\geq0$, and matches the free
Wilson response at every volume. Assume the response service chooses
$v>0,M\geq0$, depending only on these external data, with the following
property: for every sequence of positive integer sides $\ell_j\to\infty$,
every bed $B$, target $p$, path $R$, and budget $\Lambda$ satisfying the
fixed $(\tau,A,\Lambda)$ path bound, attachment of $B$ to $R$, metric
strip bound $(\sigma,C)$, direct-term bound $Z$, all grade floors
$(\epsilon,C)$, and sworn conditions at every $\ell_j$, it provides a
family $T$ with fee $(v,M)$, whose zero-, one-, and two-leg layers vanish
to orders three, two, and one respectively, and exact analytic
cancellation of the target response against the Wilson term at every
$\ell_j$. Require also the exact finite-layer realization: for every
$\ell_j$, the interacting and reference matrices are real analytic on the
real momentum torus, Hermitian positive definite at the origin with a
positive lower eigenvalue bound uniform in $j$, and their Taylor layers of
orders zero, one and two are the layers specified by the response ticket.
The resulting shifted-fiber kernel $F_{c,\ell_j}(p,q)$ is the restriction
of one common periodic kernel $F_c(p,q)$, holomorphic and bounded by some
$K$ whenever every imaginary momentum coordinate has modulus less than
some $z>0$, uniformly in $j$ and color $c$. This is the common-strip and
common-layer content needed for the full-grid limit; positivity at the
origin alone does not imply it. Assume also the uniform two-wave output: for every
combined family below, every color $c$, and every regulator $0<y\leq1$,
the same constants $a,b\geq0$ give
\[
 |\kappa_{y,c}(x)|\leq a(|x|+y)^{-4},\qquad
 |\kappa_{y,c}(x)-\Pi_{G,c}(x)|
      \leq b|x|^{-3}+by^2|x|^{-6}\quad(y\leq|x|\leq1)
 \tag{RY9}
\]
on the centered momentum box. This states the uniform local estimate
needed from the next producer; it is not deduced from the finite-volume
recording argument.

\begin{theorem}[All-period recording and drift]
\label{thm:ex-all-period-recording}
Under these hypotheses, for every admissible path datum, every odd integer
$L\geq1$, and every side sequence $\ell_j\to\infty$ there are a picture
$\mathcal Y_L$, constants
$z_L>0$, $K_L\geq0$, and numbers $e_{L,j}\geq0$ such that
\[
 \mathcal Y_L\in\mathcal H_{z_L/2}(K_L),\qquad
 e_{L,j}=\mathfrak r(z_L,K_L)e^{-z_L\ell_j/4}\longrightarrow0,
 \qquad
 \operatorname{Err}_{\ell_j}(\mathcal Y_L,\mathcal P_{\ell_j};e_{L,j}),
 \tag{RY10}
\]
and, for every color $c$,
\[
 \left|\alpha_c(\mathcal Y_L)+\lambda_G\log L\right|
 \leq D_0,
 \quad D_0=c_{\mathrm{wt}}\left[
     a\left(\frac{\pi^2}{2}+|E_4|\right)+3\pi^2 b\right].
 \tag{RY11}
\]
The constant $D_0$ is independent of $L$ and of the side sequence.
\end{theorem}
\begin{proof}
Choose the baseline before the response family. Put $w=\min(k,v)>0$;
restrict both fee bounds to $w$ and add the families, obtaining fee
$(w,M+E)$. Now fix $L$ and set
\[
 y=L^{-1},\qquad
 s=\min\{\sigma,w,y/128\}>0.
 \tag{RY5}
\]
Then $0<y\leq1$, and $s\leq y/128$ is the phase condition in the local
determinant-floor estimate. At every $\ell_j$, the response ticket matches
the target minus the Wilson term, while the baseline ticket matches the
Wilson term. Adding their zero-, one-, and two-leg Taylor identities gives
the target ticket. The finite determinant formula (RY8) turns this into
the exact local coefficient identity (RY3), and the floor and positive
regulator bounds ensure both determinants are nonzero at zero. The all-grade
floors and metric assumptions identify the regulated interacting table
with the tested local determinant table; the reference table is identified
by its explicit regulator formula. These are the two tab identities
required by the local positivity/floor argument at scale $s$. The exact
finite-layer and common-strip receipt in the hypotheses supplies
$z_L>0$, $K_L\geq0$ and the common analytic two-momentum kernel in
$\mathcal H_{z_L}(K_L)$. Applying
Lemma~\ref{lem:ex-pole-recording} yields
\[
 e_{L,j}=\mathfrak r(z_L,K_L)e^{-z_L\ell_j/4},\qquad
 \operatorname{Err}_{\ell_j}(\mathcal Y_L,\mathcal P_{\ell_j};e_{L,j}).
 \tag{RY12}
\]
The positive $z_L$ and $\ell_j\to\infty$ imply $e_{L,j}\to0$ for this
fixed $L$. No uniformity of $z_L,K_L$ in $L$ is required or asserted.

The same response satisfies (RY9), uniformly in color $c$, at regulator
$y$. The coefficient-integral identity (AT1) identifies
$\alpha_c(\mathcal Y_L)$ with the normalized integral of
$\operatorname{Re}\kappa_{y,c}$. The radial calculation in the preceding
section then gives the estimate, with the same $D_0$ for every color.
(RY11): $\log y^{-1}=\log L$, while $a,b$ were fixed uniformly before
choosing $L$. Taking $L_{\min}=1$ supplies the all-odd-side clause needed by the
outer construction; oddness is not otherwise used in this estimate.
\end{proof}

\paragraph{Scope.}
This proves the finite Fourier tail, the exact conversion of local
coefficient identities into a site-count-scaled recording error, and the
downstream assembly from the stated baseline, response-service,
local-positivity/floor, and uniform two-wave outputs. It does not reprove the local
stationary-family producers of the all-centered and shell bounds, nor the
general path-history service that supplies the response family. Those are
the remaining analytic inputs recorded in the companion evidence. No
finite-volume coefficient is identified with the continuum integral
without first passing through (RY12).
